\section{Overall Structure of the Proposed Thesis}
\label{sec:structure}

The final written thesis will be divided into the following chapters:

\begin{enumerate}
    \item \textbf{Introduction.} Introduces the problem of context saturation in flat tool-use architectures and motivates the shift toward progressive skill disclosure in autonomous agent systems. States the three research questions addressed by the study and outlines the structure of the remaining chapters.

    \item \textbf{Literature Review.} Surveys the foundational research on tool use and agentic skill architectures, trajectory-based agent evaluation, and the empirically documented effects of context window saturation on language model reasoning. Identifies the empirical gap addressed by this thesis: the absence of a controlled, reproducible comparison between flat tool use and progressive skill disclosure.

    \item \textbf{Methodology.} Describes the experimental design and all implementation details, including the two agent conditions (Condition~A: flat tool use; Condition~B: progressive disclosure via OpenClaw), the construction of the 20-skill library across three functional domains, the 40-task benchmark at three difficulty levels, the four language model backends (Gemma~4 and Qwen~3.5 via Ollama; Gemini~3.1~Pro and GPT-4o via cloud APIs), and the full set of outcome, trajectory, and efficiency metrics. Documents the Python orchestration controller and all measurement instrumentation.

    \item \textbf{Results.} Presents all collected measurements for both conditions and all model backends, organised by research question. Includes tables and figures reporting task completion rates, hallucinated invocation counts, trajectory compliance scores, prompt token consumption, end-to-end latency, and per-document inference cost. Reports results separately by task difficulty level and skill library size (10 skills and 20 skills).

    \item \textbf{Discussion.} Interprets the experimental results in relation to each of the three research questions. Analyses the conditions under which progressive disclosure provides the most substantial benefit, addresses the limitations of the study (including the use of a controlled synthetic skill library and a single hardware configuration for local model testing), and identifies directions for future research.

    \item \textbf{Conclusion.} Summarises the principal empirical findings and their practical implications for developers and organisations designing or deploying autonomous agent systems with large or growing capability sets.

    \item \textbf{References.} A complete bibliographic list of all papers, frameworks, datasets, and technical reports cited in the thesis.

    \item \textbf{Appendices.} Supporting materials including: the complete SKILL.md template used for all 20 skills; the full task benchmark with answer keys and difficulty classifications; the Python source code for the orchestration controller and metric computation scripts; raw results tables for all conditions, model backends, and repetitions; and the energy measurement configuration for each local model backend.
\end{enumerate}
